\paragraph{张量折叠的递推结果式与 Hankel 复杂度乘法律}
在命题 \ref{prop:pom-momq-tensor-kernel-construction} 给出的 Kronecker 型碰撞核之上，若把两条折叠系统作外部张量并要求碰撞矩按尺度逐点相乘，则最小递推特征式不再是抽象存在量，而受一个显式结式严格控制。

\begin{theorem}[张量折叠碰撞矩的结果式递推律]\label{thm:pom-tensor-fold-resultant-recursion}
固定整数 \(q\ge 1\)。
设两条折叠系统在矩阶 \(q\) 上的碰撞矩序列分别为
\[
\bigl(S_q^{(1)}(m)\bigr)_{m\ge 0},
\qquad
\bigl(S_q^{(2)}(m)\bigr)_{m\ge 0},
\]
并假设其张量折叠满足严格乘法律
\[
S_q^{(\otimes)}(m)=S_q^{(1)}(m)\,S_q^{(2)}(m)\qquad(m\ge 0).
\]
记 \(P_q^{(i)}(\lambda)\in\ZZ[\lambda]\) 为 \(\bigl(S_q^{(i)}(m)\bigr)_{m\ge 0}\) 的首一最小递推特征多项式，并令 \(d_i:=\deg P_q^{(i)}\)（\(i=1,2\)）。
定义
\[
\widehat P_q^{(1)}(\lambda,u):=u^{d_1}P_q^{(1)}(\lambda/u)\in\ZZ[\lambda,u],
\qquad
R_q(\lambda):=\operatorname{Res}_{u}\!\bigl(\widehat P_q^{(1)}(\lambda,u),\,P_q^{(2)}(u)\bigr)\in\ZZ[\lambda].
\]
则张量折叠序列 \(\bigl(S_q^{(\otimes)}(m)\bigr)_{m\ge 0}\) 必满足某条整系数线性递推；
若其首一最小递推特征多项式记为 \(P_q^{(\otimes)}(\lambda)\)，则
\[
P_q^{(\otimes)}(\lambda)\mid R_q(\lambda)\qquad\text{于 }\QQ[\lambda].
\]
此外，若在 \(\overline{\QQ}\) 中把
\[
P_q^{(1)}(\lambda)=\prod_{i=1}^{d_1}(\lambda-\alpha_i),
\qquad
P_q^{(2)}(\lambda)=\prod_{j=1}^{d_2}(\lambda-\beta_j)
\]
按重数分解，则有严格因式分解
\[
R_q(\lambda)=\prod_{i=1}^{d_1}\prod_{j=1}^{d_2}(\lambda-\alpha_i\beta_j).
\]
\end{theorem}
\begin{proof}
取 \(P_q^{(i)}\) 的伴随矩阵 \(C_i\in M_{d_i}(\ZZ)\)。
由于 \(P_q^{(i)}\) 首一且为最小递推特征多项式，存在向量 \(u_i,v_i\) 使
\[
S_q^{(i)}(m)=u_i^{\mathsf T}C_i^{\,m}v_i\qquad(m\ge 0),
\qquad
\chi_{C_i}(\lambda)=P_q^{(i)}(\lambda).
\]
由张量折叠的乘法律，
\[
S_q^{(\otimes)}(m)
=\bigl(u_1\otimes u_2\bigr)^{\mathsf T}
\bigl(C_1\otimes C_2\bigr)^{m}
\bigl(v_1\otimes v_2\bigr).
\]
因此 \(\bigl(S_q^{(\otimes)}(m)\bigr)\) 是矩阵 \(C_1\otimes C_2\) 的线性读出；于是任一满足
\(
F(C_1\otimes C_2)=0
\)
的多项式 \(F\in\QQ[\lambda]\) 都诱导一条递推
\(
F(E)S_q^{(\otimes)}\equiv 0
\)。

在 \(\overline{\QQ}\) 中，\(C_1\) 与 \(C_2\) 的特征值多重集分别为 \(\{\alpha_i\}_{i=1}^{d_1}\) 与 \(\{\beta_j\}_{j=1}^{d_2}\)，故
\[
\chi_{C_1\otimes C_2}(\lambda)
=\prod_{i=1}^{d_1}\prod_{j=1}^{d_2}(\lambda-\alpha_i\beta_j).
\]
另一方面，
\[
\widehat P_q^{(1)}(\lambda,u)
=u^{d_1}P_q^{(1)}(\lambda/u)
=\prod_{i=1}^{d_1}(\lambda-\alpha_i u),
\]
因而结式恒等式给出
\[
R_q(\lambda)
=\operatorname{Res}_{u}\!\bigl(\widehat P_q^{(1)}(\lambda,u),\,P_q^{(2)}(u)\bigr)
=\prod_{j=1}^{d_2}\widehat P_q^{(1)}(\lambda,\beta_j)
=\prod_{i=1}^{d_1}\prod_{j=1}^{d_2}(\lambda-\alpha_i\beta_j).
\]
故 \(R_q(\lambda)=\chi_{C_1\otimes C_2}(\lambda)\)。
由 Cayley--Hamilton 定理，
\(
R_q(C_1\otimes C_2)=0
\)，从而 \(R_q(E)S_q^{(\otimes)}\equiv 0\)。
由于 \(P_q^{(\otimes)}\) 按定义为该序列的首一最小递推特征多项式，必有
\(
P_q^{(\otimes)}\mid R_q
\)
于 \(\QQ[\lambda]\)。
证毕。
\end{proof}

\begin{corollary}[非碰撞制度下的 Hankel 复杂度严格乘法性]\label{cor:pom-tensor-fold-hankel-rank-multiplicativity}
沿用定理 \ref{thm:pom-tensor-fold-resultant-recursion} 的设定，并记
\(
\operatorname{rank}_H(\cdot)
\)
为序列的 Hankel 秩（等价地，其最小递推阶数）。
若进一步满足：
\begin{enumerate}
  \item \(P_q^{(1)}\) 与 \(P_q^{(2)}\) 在 \(\overline{\QQ}\) 上均无重根；
  \item 按其根展开的谱系数均非零，即存在
  \[
  S_q^{(1)}(m)=\sum_{i=1}^{d_1}a_i\alpha_i^m,
  \qquad
  S_q^{(2)}(m)=\sum_{j=1}^{d_2}b_j\beta_j^m
  \qquad(a_i,b_j\neq 0) ;
  \]
  \item 谱乘积无碰撞：
  \[
  \alpha_i\beta_j=\alpha_{i'}\beta_{j'}
  \Longrightarrow
  (i,j)=(i',j').
  \]
\end{enumerate}
则
\[
\deg P_q^{(\otimes)}=d_1d_2,
\qquad
\operatorname{rank}_H\!\bigl(S_q^{(\otimes)}\bigr)
=\operatorname{rank}_H\!\bigl(S_q^{(1)}\bigr)\,
\operatorname{rank}_H\!\bigl(S_q^{(2)}\bigr).
\]
\end{corollary}
\begin{proof}
由简单根分解与系数非零性，
\[
S_q^{(\otimes)}(m)
=S_q^{(1)}(m)S_q^{(2)}(m)
=\sum_{i=1}^{d_1}\sum_{j=1}^{d_2}(a_ib_j)\,(\alpha_i\beta_j)^m.
\]
由谱乘积无碰撞，\(\{\alpha_i\beta_j\}_{i,j}\) 两两不同。
对特征零域上的两两不同标量 \(\gamma_1,\dots,\gamma_N\)，指数序列 \(\gamma_k^m\)（\(0\le m\le N-1\)）线性无关；否则 Vandermonde 行列式
\[
\det(\gamma_k^{\,m})_{0\le m\le N-1,\ 1\le k\le N}
=\prod_{1\le k<\ell\le N}(\gamma_\ell-\gamma_k)
\]
将为零，矛盾。
因此 \(\bigl(S_q^{(\otimes)}(m)\bigr)\) 的最小递推阶数至少为 \(d_1d_2\)。

另一方面，定理 \ref{thm:pom-tensor-fold-resultant-recursion} 表明
\(
P_q^{(\otimes)}\mid R_q
\)，而 \(\deg R_q=d_1d_2\)。
故 \(\deg P_q^{(\otimes)}\le d_1d_2\)。
两侧合并即得 \(\deg P_q^{(\otimes)}=d_1d_2\)。
再由 Hankel 秩与最小递推阶数的一致性，即得所示严格乘法律。证毕。
\end{proof}

\endinput
